\documentclass[11pt]{amsart}
\usepackage{amsmath,amssymb,amsthm}
\usepackage[margin=1.05in]{geometry}
\raggedbottom

\numberwithin{equation}{section}
\newtheorem{theorem}{Theorem}[section]
\newtheorem{proposition}[theorem]{Proposition}
\newtheorem{lemma}[theorem]{Lemma}
\newtheorem{corollary}[theorem]{Corollary}
\theoremstyle{definition}
\newtheorem{definition}[theorem]{Definition}
\theoremstyle{remark}
\newtheorem{remark}[theorem]{Remark}
\newtheorem{example}[theorem]{Example}

\newcommand{\Z}{\mathbb{Z}}
\newcommand{\Q}{\mathbb{Q}}
\newcommand{\R}{\mathbb{R}}
\newcommand{\N}{\mathbb{N}}
\newcommand{\C}{\mathbb{C}}
\newcommand{\cA}{\mathcal{A}}
\newcommand{\cG}{\mathcal{G}}
\newcommand{\cO}{\mathcal{O}}
\newcommand{\cL}{\mathcal{L}}
\newcommand{\pp}{\mathfrak{p}}
\newcommand{\aaa}{\mathfrak{a}}
\newcommand{\Ns}{\mathrm{N}}
\newcommand{\Cl}{\operatorname{Cl}}
\newcommand{\Gal}{\operatorname{Gal}}
\newcommand{\Iso}{\operatorname{Iso}}
\newcommand{\Idl}{\mathbb{I}}
\newcommand{\rank}{\operatorname{rank}}
\newcommand{\id}{\mathrm{id}}

\PassOptionsToPackage{hyphens}{url}
\usepackage[hidelinks,bookmarks=false]{hyperref}

\title[Leopoldt, the strata, and $\log_p(\Ns\pp)=0$]{Leopoldt and the $p$-local Bost--Connes groupoid:\\ the rank of the isotropy, not its triviality,\\ and why $\log_p(\Ns\pp)$ vanishes}
\author{Ruqing Chen}
\address{GUT Geoservice Inc., Montreal, Canada}
\email{ruqing@hotmail.com}
\thanks{Sources, code and data for the entire series: \url{https://github.com/Ruqing1963/localized-bost-connes}; this paper is in \texttt{papers/XXIII-leopoldt/}.}
\date{}

\begin{document}

\begin{abstract}
Paper XXII of the series identified two genuine points of contact between the localized
Bost--Connes systems and $p$-adic arithmetic. We develop both and find that each of the three
theorems proposed on that basis needs correction, in each case to a statement that is true and
cleaner.

\emph{Leopoldt does not make the intermediate isotropy trivial --- and the isotropy is not
the group the proposal computed.} By Proposition 2.2 of Paper II of the series, the isotropy
at the stratum $Y_T$ is the \emph{divisor image}
$\mathrm{div}\bigl(\overline{K^\times_{S,+}}\cap\Idl_T\bigr)$, and we identify it exactly:
$\Iso(Y_T)\cong\ker\bigl(\varphi:L_T\to U_{T^c}/\pi_{T^c}(\overline{\cO^\times_{K,+}})\bigr)$,
where $L_T$ is the rank-$|T|$ lattice of $T$-supported principal divisors with totally
positive generator and $\pi_{T^c}$ is the projection of the closed unit group to the
complementary factors. The group $\overline{\cO^\times_{K,+}}\cap\prod_{\pp\in T}\cO_\pp^\times$
that the proposal analyzed is the \emph{kernel} of the divisor map --- it contributes the
identity ideal --- though its rank, $(r-\delta_p(K))-\rank\pi_{T^c}$, is exactly what budgets
$\rank\pi_{T^c}$. The consequence is stronger than the proposed correction: whenever
$d_{T^c}=\sum_{\pp\notin T}[K_\pp:\Q_p]=1$ and $r\ge1$, the image of $\pi_{T^c}$ is open ---
a single coordinate of an infinite-order unit has nonzero logarithm --- so the quotient is
finite and $\Iso(Y_T)$ has \emph{full rank} $|T|$. Already for a real quadratic field with
$p$ split the intermediate isotropy is nontrivial, unconditionally. The proposed equivalence
``Leopoldt $\iff$ trivial isotropy'' is therefore false; Leopoldt enters only as the budget
$\rank\pi_{T^c}\le r-\delta_p(K)$.

\emph{The proposed $\cL$-invariant is identically zero.} The derivative computation is
correct: $E_p(\beta)=\prod_{\pp\mid p}(1-\Ns\pp^{-\beta})$ vanishes to order $g=\#\{\pp\mid p\}$
at $\beta=0$ with
$\frac{1}{g!}E_p^{(g)}(0)=\prod_{\pp\mid p}\log\Ns\pp=(\log p)^g\prod_\pp f_\pp$. But this is
the \emph{real} logarithm. Its $p$-adic counterpart vanishes: Iwasawa's branch has
$\log_p(p)=0$, so $\log_p(\Ns\pp)=f_\pp\log_p(p)=0$ for every $\pp\mid p$ and the proposed
$\cL_p(K)=\prod_\pp\log_p(\Ns\pp)$ is $0$. It cannot be a Fontaine--Mazur invariant.

\emph{The regulator is a dimension, not a volume.} On a compact group Haar measure never
degenerates, and $\overline{\cO^\times_{K,+}}$ is infinite, so the quotient has no
informative volume. What is meaningful is
\[
\dim\Bigl(U_p\big/\overline{\cO^\times_{K,+}}\Bigr)=[K:\Q]-\bigl(r-\delta_p(K)\bigr)=r_2+1+\delta_p(K),
\]
so that Leopoldt is equivalent to this dimension being exactly $r_2+1$ --- recovering
Proposition 5.3 of Paper XXII.

What survives intact is the order of vanishing: $E_p$ vanishes at $\beta=0$ to order exactly
$g$, and $g$ is the number of exceptional zeros in the Mazur--Tate--Teitelbaum sense. That
part of the analogy is correct and is the one worth keeping.
\end{abstract}

\maketitle

\section{The isotropy of the intermediate strata}

Let $K$ be a number field, $p$ a rational prime, $S=\{\pp\mid p\}$ with $g=\#S$, and
$U_p=\prod_{\pp\mid p}\cO_\pp^\times$. Write $r=r_1+r_2-1$ for the unit rank, $\delta_p(K)\ge0$
for the Leopoldt defect, so that $\rank_{\Z_p}\overline{\cO^\times_{K,+}}=r-\delta_p(K)$
\cite{Washington} --- $\delta_p(K)=0$ is known for abelian $K$ by Baker--Brumer
\cite{Brumer} --- and
for $T\subseteq S$ put $d_{T^c}=\sum_{\pp\notin T}[K_\pp:\Q_p]$ and
$\pi_{T^c}:\overline{\cO^\times_{K,+}}\to\prod_{\pp\notin T}\cO_\pp^\times$ for the
coordinate projection. The framework is that of \cite{Main}; Papers II and XXII of the
series, referred to by numeral, are unpublished companions, and every fact used from them is
restated or verified in place.

\begin{proposition}[Kernel and projection ranks]\label{prop:rank}
The group $\overline{\cO^\times_{K,+}}\cap\prod_{\pp\in T}\cO_\pp^\times=\ker\pi_{T^c}$
satisfies
\begin{equation}\label{eq:rank}
\rank_{\Z_p}\ker\pi_{T^c}
=\bigl(r-\delta_p(K)\bigr)-\rank\pi_{T^c}
\ \ge\ \max\bigl(0,\ r-\delta_p(K)-d_{T^c}\bigr),
\end{equation}
with equality if and only if $\pi_{T^c}$ has maximal rank $\min(r-\delta_p,d_{T^c})$.
Moreover $\rank\pi_{T^c}\ge1$ whenever $r-\delta_p\ge1$ and $T^c\ne\emptyset$: a coordinate
of an infinite-order unit is an infinite-order element of $\cO_{\pp'}^\times$, its logarithm
being nonzero since $K\hookrightarrow K_{\pp'}$ is injective. We emphasize that this kernel
is \emph{not} the isotropy of the stratum: it is the divisor-zero part, mapping to the
identity ideal.
\end{proposition}

\begin{proof}
The $p$-adic logarithm embeds $\overline{\cO^\times_{K,+}}$, up to finite index, as a
$\Z_p$-lattice of rank $r-\delta_p$ in $\bigoplus_{\pp\mid p}K_\pp$, of total dimension
$[K:\Q]$; the kernel of the coordinate projection onto
$\bigoplus_{\pp\notin T}K_\pp$, of dimension $d_{T^c}$, has rank given by rank-nullity,
which is \eqref{eq:rank}.
\end{proof}

\begin{theorem}[The isotropy, identified; Leopoldt does not give triviality]\label{thm:nottrivial}
Let $T\subsetneq S$ and let $L_T\subseteq\Z^T$ be the finite-index sublattice of exponent
vectors $a$ with $\prod_{\pp\in T}\pp^{a_\pp}$ principal with totally positive generator
$\alpha_a$. By Proposition 2.2 of Paper II,
$\Iso(y)=\mathrm{div}\bigl(\overline{K^\times_{S,+}}\cap\Idl_T\bigr)$ for $y\in Y_T$, and
this equals, exactly,
\[
\Iso(Y_T)\ \cong\ \ker\Bigl(\varphi:L_T\longrightarrow
\prod_{\pp\notin T}\cO_\pp^\times\Big/\pi_{T^c}\bigl(\overline{\cO^\times_{K,+}}\bigr)\Bigr),
\qquad \varphi(a)=\alpha_a\big|_{T^c},
\]
well defined since a change of generator moves $\alpha_a|_{T^c}$ by
$\pi_{T^c}(\cO^\times_{K,+})$, and with $\pi_{T^c}(\overline{\cO^\times_{K,+}})$ closed, the
closed unit group being compact. Consequently:
\begin{enumerate}
\item if $\cO^\times_{K,+}$ is finite, $\varphi(a)$ torsion forces $\alpha_a$ to equal a
root of unity in some completion, hence in $K$, so $a=0$: the isotropy is trivial ---
Paper II's unconditional case;
\item if $r\ge1$ and $d_{T^c}=1$, the image of $\pi_{T^c}$ is \emph{open} in the
one-dimensional $\cO_{\pp'}^\times$ (Proposition~\ref{prop:rank}, last part: a nonzero
$\Z_p$-multiple of a nonzero logarithm is open), the quotient is finite, and
$\Iso(Y_T)$ has \emph{full rank} $|T|$. In particular, for $K$ real quadratic with $p$
split and $T$ one of the two primes, the isotropy is infinite --- \emph{unconditionally};
for $K$ totally real of degree $n$ with $p$ split completely and $|T^c|=1$ it has rank
$n-1$.
\end{enumerate}
The proposed equivalence
\[
\text{Leopoldt}\iff\Iso(y)=\{1\}\ \text{for all}\ T\subsetneq S
\]
is therefore false: triviality tracks the finiteness of the unit group, not Leopoldt, and
fails under Leopoldt as soon as $r\ge1$ (with a small complementary block).
\end{theorem}

\begin{proof}
For $a\in\ker\varphi$ pick $\beta\in\overline{\cO^\times_{K,+}}$ with
$\pi_{T^c}(\beta)=\alpha_a|_{T^c}$; then $\alpha_a\beta^{-1}\in
\overline{K^\times_{S,+}}\cap\Idl_T$ has divisor $a$, so $a\in\Iso$. Conversely if
$\gamma\in\overline{K^\times_{S,+}}\cap\Idl_T$ has divisor $a$, approximate $\gamma$ by
totally positive $S$-units $\gamma_n$; the divisor map is locally constant, so
$v(\gamma_n)=a$ eventually, whence $\gamma_n=\alpha_au_n$ with
$u_n\in\cO^\times_{K,+}$ (times torsion), and $\pi_{T^c}(u_n)\to\alpha_a|_{T^c}^{-1}$; the
image $\pi_{T^c}(\overline{\cO^\times_{K,+}})$ is compact, hence closed, so
$\alpha_a|_{T^c}^{-1}$ lies in it and $\varphi(a)=0$. Items (1) and (2) follow as indicated;
for the rank in (2), $L_T$ has rank $|T|$ and the quotient is finite, so
$\ker\varphi$ has finite index in $L_T$.
\end{proof}

\begin{remark}[Numerical table]\label{rem:table}
With $p$ split completely and $|T^c|=1$, so $|T|=g-1=n-1$ and $d_{T^c}=1$, the corrected
isotropy rank is:
\[
\begin{array}{lcccc}
K & r & |T| & \text{kernel rank \eqref{eq:rank}} & \rank\Iso(Y_T)\\\hline
\Q & 0 & 0 & 0 & 0\\
\text{imaginary quadratic} & 0 & 1 & 0 & 0\\
\text{real quadratic} & 1 & 1 & 0 & 1\\
\text{real cubic} & 2 & 2 & 1 & 2\\
\text{real quartic} & 3 & 3 & 2 & 3\\
\text{real quintic} & 4 & 4 & 3 & 4
\end{array}
\]
The last column is Theorem~\ref{thm:nottrivial}(2) (full rank $|T|$ as soon as $r\ge1$),
unconditional; the proposal's values were the kernel column, which is not the isotropy.
Triviality holds exactly in the first two rows --- the finite-unit-group cases, which are
precisely those in which Paper II (its Proposition 2.2 and Remark 2.3) is unconditional.
\end{remark}

\begin{remark}[Nothing in the classification depends on this]\label{rem:nodependence}
The isotropy on the strata does not enter the main results. Paper II
established topological principality of $\cG_S$ \emph{directly} (its Theorem 2.5), the
strata being nowhere
dense regardless of their isotropy, which is what Renault's machinery requires. So
Theorem~\ref{thm:nottrivial} corrects a proposed characterization of Leopoldt without
disturbing anything downstream.
\end{remark}

\begin{remark}[The residual condition, and where Leopoldt sits]\label{rem:relative}
For $d_{T^c}\ge2$ the finiteness of the quotient in Theorem~\ref{thm:nottrivial} requires
$\pi_{T^c}$ to have \emph{open} image, i.e.\ rank $d_{T^c}$ --- a relative nondegeneracy of
the regulator matrix restricted to the columns of $T^c$, automatic for $d_{T^c}=1$ but in
general neither implied by nor implying Leopoldt, which is the single statement that the
full $r\times r$ $p$-adic regulator is nonzero. Leopoldt enters only as the budget: all
projection ranks are at most $r-\delta_p(K)$, so a positive defect can only \emph{lower}
$\rank\pi_{T^c}$ and thereby raise the codimension of the image.
\end{remark}

\section{The exceptional Euler factor}

\begin{proposition}[Order of vanishing and leading term]\label{prop:euler}
Let $E_p(\beta)=\zeta_{K,p}(\beta)^{-1}=\prod_{\pp\mid p}\bigl(1-\Ns\pp^{-\beta}\bigr)$. Then
$E_p$ vanishes at $\beta=0$ to order exactly $g=\#\{\pp\mid p\}$, and
\begin{equation}\label{eq:leading}
\frac{1}{g!}\,\frac{d^g}{d\beta^g}E_p(\beta)\Big|_{\beta=0}
=\prod_{\pp\mid p}\log\Ns\pp=(\log p)^g\prod_{\pp\mid p}f_\pp ,
\end{equation}
$f_\pp$ being the residue degrees.
\end{proposition}

\begin{proof}
$1-\Ns\pp^{-\beta}=1-e^{-\beta\log\Ns\pp}=\beta\log\Ns\pp+O(\beta^2)$, and
$\Ns\pp=p^{f_\pp}$.
\end{proof}

\begin{theorem}[The proposed $\cL$-invariant vanishes]\label{thm:Lzero}
With Iwasawa's branch of the $p$-adic logarithm, normalized by $\log_p(p)=0$, one has
\[
\log_p(\Ns\pp)=\log_p\bigl(p^{f_\pp}\bigr)=f_\pp\log_p(p)=0
\qquad\text{for every }\pp\mid p ,
\]
so $\prod_{\pp\mid p}\log_p(\Ns\pp)=0$. The quantity proposed as
$\cL_p(K)$ is identically zero and cannot be a Fontaine--Mazur or
Mazur--Tate--Teitelbaum $\cL$-invariant \cite{MTT,GreenbergSteven}: those are nonzero in
the proved cases --- for split multiplicative elliptic curves
$\cL_p=\log_p(q_E)/\mathrm{ord}_p(q_E)\ne0$ --- and conjecturally always.
\end{theorem}

\begin{remark}[Real versus $p$-adic]\label{rem:realvsp}
The confusion is between two logarithms. Formula \eqref{eq:leading} is a computation with the
\emph{real} logarithm, coming from the archimedean dynamics
$\sigma_t(\mu_\aaa)=\Ns\aaa^{it}\mu_\aaa$, and it is correct; it evaluates to
$(\log p)^g\prod f_\pp$, for instance $2.5903$ for $g=2$, $f=(1,1)$, $p=5$, and $4.1689$ for
$g=3$, $f=(1,1,1)$, $p=5$. The $p$-adic logarithm of the same quantity is $0$, because
$\Ns\pp$ is a power of $p$ and $\log_p$ is normalized to kill $p$. The two cannot be
identified, and the vanishing is not an accident: it is the same fact as
Proposition 2.1 of Paper XXII, that $\Ns\pp$ is not a $p$-adic unit.
\end{remark}

\begin{remark}[What is right about the analogy]\label{rem:rightpart}
The order of vanishing is correct and is worth stating on its own. $E_p$ vanishes to order
$g$ at $\beta=0$, and $g$ is exactly the number of primes above $p$ at which the Euler factor
of $L_p(s,\chi)$ is removed, hence the number of exceptional zeros in the
Mazur--Tate--Teitelbaum sense \cite{MTT} for the trivial character. That is a genuine dictionary entry,
and it is the one supported by Proposition 5.1 of Paper XXII: the $p$-local partition function
corresponds to the interpolation factor. The leading coefficient of the archimedean $E_p$ is
not the $\cL$-invariant of that factor.
\end{remark}

\section{The regulator is a dimension, not a volume}

\begin{proposition}[The volume statement is not well formed]\label{prop:novolume}
$G_p$ is a compact group, so its Haar measure exists, is unique up to scale and is never
degenerate; normalizing $m_{G_p}(G_p)=1$ makes any statement about its total mass vacuous.
Moreover $\overline{\cO^\times_{K,+}}$ is infinite whenever $r>\delta_p(K)$, so the index
$[U_p:\overline{\cO^\times_{K,+}}]$ is infinite and the naive relative volume is $0$
independently of $R_p(K)$. No equivalence with $R_p(K)\ne0$ can be extracted from a volume.
\end{proposition}

\begin{theorem}[The correct statement]\label{thm:dimension}
\[
\dim\Bigl(U_p\big/\overline{\cO^\times_{K,+}}\Bigr)
=[K:\Q]-\bigl(r-\delta_p(K)\bigr)
=r_2+1+\delta_p(K),
\]
as $p$-adic analytic groups. Consequently
\[
\text{Leopoldt for }K\text{ at }p\iff R_p(K)\ne0\iff \dim\Bigl(U_p\big/\overline{\cO^\times_{K,+}}\Bigr)=r_2+1 .
\]
\end{theorem}

\begin{proof}
$\dim U_p=\sum_{\pp\mid p}[K_\pp:\Q_p]=[K:\Q]=r_1+2r_2$ and
$\dim\overline{\cO^\times_{K,+}}=r-\delta_p$ with $r=r_1+r_2-1$; subtract. The equivalence is
the definition of the Leopoldt defect together with the fact that $\delta_p=0$ iff the
$p$-adic regulator is nonzero.
\end{proof}

\begin{example}\label{ex:dim}
Under Leopoldt: $\dim=1$ for $\Q$, for real quadratic, for real cubic and for real quartic
fields; $\dim=2$ for imaginary quadratic and for complex cubic fields. In each case
$\dim=r_2+1$.
\end{example}

\begin{corollary}[Agreement with Paper XXII]\label{cor:agree}
Theorem~\ref{thm:dimension} is the computation
\[
\rank_{\Z_p}\Gal(M_\infty/K)=r_2+1+\delta_p(K)
\]
of Paper XXII (its Proposition 5.3), transported through class field theory:
$\Gal(M_\infty/K)$ is the quotient of $U_p$ by the
closure of the global units, up to a finite group. The intended statement was therefore
correct in substance and misstated in form, volume for dimension.
\end{corollary}

\section{Assessment}

\[
\begin{array}{lll}
\text{Proposal} & \text{Status} & \text{Reference}\\\hline
\Iso(Y_T)\cong\ker\varphi\ \text{(divisor image)} & \textbf{exact identification} & \text{Thm.~1.2}\\
\text{full rank }|T|\ \text{for }r\ge1,\ d_{T^c}=1 & \textbf{unconditional} & \text{Thm.~1.2}\\
\text{Leopoldt}\iff\text{trivial isotropy} & \textbf{false} & \text{Thm.~1.2}\\
\text{topological principality affected} & \textbf{no} & \text{Rem.~1.4}\\
\mathrm{ord}_{\beta=0}E_p=g & \text{true} & \text{Prop.~2.1}\\
\tfrac{1}{g!}E_p^{(g)}(0)=\prod\log\Ns\pp & \text{true, real logarithm} & \text{Prop.~2.1}\\
\cL_p(K)=\prod\log_p(\Ns\pp) & \textbf{identically }0 & \text{Thm.~2.2}\\
g=\#\{\text{exceptional zeros}\} & \textbf{true; keep this} & \text{Rem.~2.4}\\
\mathrm{Vol}_p(U_p/\overline{\cO^\times})\asymp|R_p|_p & \textbf{not well formed} & \text{Prop.~3.1}\\
\dim=r_2+1+\delta_p & \textbf{true; the right form} & \text{Thm.~3.2}
\end{array}
\]

The two bridges identified in Paper XXII are real, and both survive; what fails is the
attempt to make each carry a theorem it cannot. On the unit side, Leopoldt controls the
\emph{rank} of the closed unit group and hence the \emph{dimension} of the symmetry group,
and that is the whole of its role: it does not trivialize the isotropy of the strata, and the
strata do not matter for the classification in any case. On the Euler side, the $p$-local
partition function does correspond to the interpolation factor, and the order of vanishing
does count the exceptional zeros; but the leading coefficient is an archimedean quantity, and
its $p$-adic shadow is zero because $\Ns\pp$ is a power of $p$.

That last point is worth isolating, because it recurs. In Proposition 2.1 of Paper XXII the
$p$-adic dynamics failed to exist at $\pp\mid p$ because $\Ns\pp\notin\Z_p^\times$; here the
$\cL$-invariant vanishes for the same reason. Every attempt to run $p$-adic analysis on the
norms of the primes above $p$ meets the same wall, and it is not a technical one: those norms
are precisely the elements the $p$-adic theory is built to discard.

Two questions. Is there a version of the isotropy statement that does characterize Leopoldt
--- one would want a condition on \emph{all} $T$ simultaneously, and since by
Proposition~\ref{prop:rank} every projection rank is budgeted by $r-\delta_p(K)$, some
aggregate of the image codimensions does see $\delta_p$ and might work. And, given Remark~\ref{rem:rightpart}, is the correspondence
between the order of vanishing of $E_p$ and the exceptional-zero count functorial enough to
extend to nontrivial characters $\chi$, where the order becomes
$\#\{\pp\mid p:\chi(\pp)=1\}$ rather than $g$?

\begin{thebibliography}{9}
\bibitem{Brumer} A. Brumer, \emph{On the units of algebraic number fields}, Mathematika 14 (1967), 121--124.
\bibitem{GreenbergSteven} R. Greenberg, G. Stevens, \emph{$p$-adic $L$-functions and $p$-adic periods of modular forms}, Invent. Math. 111 (1993), 407--447.
\bibitem{MTT} B. Mazur, J. Tate, J. Teitelbaum, \emph{On $p$-adic analogues of the conjectures of Birch and Swinnerton-Dyer}, Invent. Math. 84 (1986), 1--48.
\bibitem{Washington} L. C. Washington, \emph{Introduction to Cyclotomic Fields}, 2nd ed., Graduate Texts in Math. 83, Springer, 1997.
\bibitem{Main} R. Chen, \emph{KMS state uniqueness and phase transitions in localized Bost--Connes systems}, preprint, 2026, \newline\url{https://doi.org/10.5281/zenodo.22152101}.
\end{thebibliography}

\end{document}
